\section{Discussion}
\label{sec:discussion}

The benchmark in Section~\ref{sec:benchmark} produced two headline findings that run against the standard intuition for shrinkage estimation: (i)~the best standalone precision estimator (direct precision shrinkage with a matched target) is not the best sequential estimator (Ledoit--Wolf), and (ii)~direct precision shrinkage with a diagonal target fails spectacularly on a system with dense off-diagonal precision (Lorenz~96) even though the framework is formally correct. This section develops the explanation.

\subsection{Target-Truth Alignment}
\label{sec:disc_alignment}

The reversal between Lorenz~96 and cosmological data is driven by whether the target $\boldsymbol{\Pi}_0$ captures the dominant structure of the true precision. On BOSS-like data the true precision is near-diagonal by construction, and any of the diagonal targets we tested (identity, eigenvalue, cosmological, scaled identity) are a reasonable match; the Bodnar coefficients stay in $[0,1]$, very little clamping occurs, and the method performs as the theory predicts. On Lorenz~96 the true precision is \emph{dense} because nearest-neighbor coupling produces pervasive off-diagonal structure, yet all four targets are diagonal by construction. The Cauchy--Schwarz denominator in Eq.~\eqref{eq:alpha_star} becomes small relative to the numerator, $\beta^*$ explodes to $10^8$--$10^9$ for the identity target and $10^3$--$10^5$ for the scaled variant, and clamping to $[0,1]$ fires in over 80\% of cycles. The resulting precision matrix is a clamped blend of two mismatched ingredients, with condition numbers of $10^7$--$10^{11}$, and the EnKF update amplifies the error.

The eigenvalue target is a diagnostic counter-example: because it is derived from the same sample as $\hat{\boldsymbol{\Pi}}_S$, the Cauchy--Schwarz term $\|\hat{\boldsymbol{\Pi}}_S\|_F^2\,\|\boldsymbol{\Pi}_0^{(1)}\|_F^2 - (\mathrm{tr}(\hat{\boldsymbol{\Pi}}_S \boldsymbol{\Pi}_0^{(1)}))^2$ is tiny and drives $\beta^* \approx 10^{-4}$. The estimator collapses back to $\alpha^* \hat{\boldsymbol{\Pi}}_S$---a rescaled sample precision with no effective shrinkage---so the eigenvalue target's poor Lorenz~96 performance is a symptom of self-reference, not of the framework.

\subsection{Variance, not Bias, Dominates in Sequential Filtering}
\label{sec:disc_variance}

The more surprising observation is that Ledoit--Wolf wins the cosmological EnKF benchmark (Table~\ref{tab:cosmo_enkf}) despite losing the standalone Frobenius-loss benchmark by a factor of two (Table~\ref{tab:cosmo_direct}). This is not an artefact: the standalone experiment gives each estimator one sample; the filtering experiment gives it one per cycle for 100 cycles, and the analysis error at step $t$ feeds the forecast error at step $t+1$ through the model dynamics.

A simple accumulation argument makes this quantitative. If an estimator has per-step bias $b$ and per-step variance $s^2$, the total error squared over $K$ cycles scales roughly as $K\,b^2 + K\,s^2$ for approximately independent cycles. Ledoit--Wolf's standard deviation across ensemble draws is $\pm 0.01$, an order of magnitude below the direct precision methods' $\pm 0.13$--$0.14$. A $100\times$ reduction in $s^2$ easily outweighs a $2\times$ penalty in $b^2$ once $K$ is in the hundreds. The bias--variance tradeoff does not reverse between settings; the sequential setting merely re-weights it by the number of cycles.

This observation points to a gap in the shrinkage-estimation literature. Ledoit and Wolf~\cite{ledoit2004well}, Chen et al.~\cite{chen2010shrinkage}, and Bodnar et al.~\cite{bodnar2016direct} all optimise single-shot loss. Review articles of the EnKF literature describe the sequential structure but do not connect it to estimator choice; the consequences are typically absorbed into inflation tuning rather than estimator selection. Our results suggest that, for sequential use, the \emph{stability} of the estimator (its variance across realisations) is at least as important as its accuracy, and that this criterion is not visible in a standalone comparison.

\subsection{Practical Recommendations}
\label{sec:disc_recommendations}

Combining the accuracy and cost results, a small decision table emerges. When $N_e \leq n$ (rank-deficient), only localization and the modified Cholesky variant apply; direct precision shrinkage is undefined. When $N_e \approx n$ (ill-conditioned) with near-diagonal precision and a known domain-specific target, direct precision shrinkage with that target gives the best standalone estimate; for sequential use, Ledoit--Wolf is the safer choice because of its low cycle-to-cycle variance. When the precision is dense, Ledoit--Wolf (or EnKF-Cholesky at intermediate $n$) is preferred; diagonal direct-precision targets fail. When $N_e \gg n$ (well-sampled), the sample precision is already accurate, and all shrinkage methods converge toward it. In all cases, NERCOME is ruled out for sequential use by its $\mathcal{O}(B)$ eigendecomposition cost.

The scaled identity target we propose occupies the niche where the precision is near-diagonal but no domain-specific target is available. It gives modest gains over the plain identity without requiring external information, and matches or beats Ledoit--Wolf on standalone estimation at small $n$. Its ad-hoc factor of~2, borrowed by analogy from the cosmological $2/N_\ell$ term, is admittedly heuristic and could likely be improved through cross-validation or a Bayesian argument.

\subsection{Limitations and Future Work}
\label{sec:disc_future}

Three limitations matter for interpreting the results. First, all experiments use a linear observation operator; nonlinear operators would introduce additional approximation error that can interact with the precision estimation error. Second, both test systems are low-dimensional ($n = 40$ and $d = 18$) compared to operational data assimilation ($n \sim 10^6$--$10^8$); the cost extrapolation in Section~\ref{sec:results} implies that the $\mathcal{O}(n^3)$ inversion becomes prohibitive at $n \gtrsim 10^4$ unless the direct-precision framework is combined with localization or a low-rank approximation. Third, the sequential-accumulation effect is empirical; a rigorous Lyapunov-style analysis of how estimator variance propagates across filtering cycles would put our finding on firmer theoretical ground.

These limitations suggest three concrete extensions. A hybrid that first localises the sample covariance and then applies direct precision shrinkage could carry the framework to $n \sim 10^3$--$10^4$, though the localized matrix is no longer Wishart and the Bodnar derivation would need revisiting. An adaptive target that switches between diagonal and covariance-space methods based on a cheap off-diagonal-energy diagnostic would handle both Lorenz-like and cosmology-like regimes within a single estimator. And a condition-number-capped variant of the Bodnar estimator, truncating tail eigenvalues before coefficient computation, would directly address the failure mode observed on Lorenz~96.
